\documentclass[11pt]{article}
\usepackage[margin=1in]{geometry}
\usepackage{amsmath,amssymb,amsthm}
\usepackage{microtype}
\usepackage[hidelinks]{hyperref}

\newtheorem{theorem}{Theorem}
\newtheorem{proposition}{Proposition}
\newcommand{\SU}{\mathrm{SU}(2)}
\newcommand{\cL}{\mathcal L}
\newcommand{\HS}{\mathrm{HS}}
\newcommand{\N}{\mathbb N}

\title{A Dirac-mass counterexample to the printed\\
Beurling ultradistribution criterion on \(\mathrm{SU}(2)\)}
\author{A full negative answer and corrected theorem for Conjecture 4.5.5\\
in arXiv:1804.03544}
\date{11 August 2026}

\begin{document}
\maketitle

\noindent\textbf{Status.}
\texttt{candidate\_full\_counterexample\_likely\_valid\_needs\_human\_review}.

\begin{abstract}
The second clause of Conjecture 4.5.5 in the source characterizes the dual
of a sub-Laplacian Gevrey--Beurling class on \(\SU\) using a \emph{positive}
exponential Fourier weight.  It is false.  The Dirac mass at the identity is
an ordinary distribution and hence belongs to that dual, as the source itself
records, but its Fourier blocks are identity matrices.  The proposed norms
grow at least as \(\sqrt2 e^{B l^{1/(2s)}}\).  We also prove the exact repair:
the exponent in the second clause must be negative.  With that correction,
the standard weighted Hilbert-scale argument proves both clauses.
\end{abstract}

\section{The printed conjecture}

Let \(X,Y\) be the canonical H\"ormander pair on \(\SU\) and
\(\cL=-(X^2+Y^2)\).  Irreducible representations are indexed by
\(l\in\frac12\N_0\), have dimension \(d_l=2l+1\), and the source gives
\begin{equation}\label{eq:symbol}
 \widehat{\cL}(l)_{m,n}
   =\bigl(l(l+1)-m^2\bigr)\delta_{m,n},
 \qquad -l\le m,n\le l.
\end{equation}
For \(1\le s<\infty\), put \(a=1/(2s)\).  The first clause of source
Conjecture 4.5.5 asserts that the Roumieu dual consists exactly of those
\(u\) for which, for every \(B>0\),
\begin{equation}\label{eq:roumieu-source}
 \sup_l\|e^{-B\widehat{\cL}(l)^a}\widehat u(l)\|_{\HS}<\infty.
\end{equation}
The second clause asserts that \(v\) belongs to the Beurling dual exactly
when, for some \(B>0\),
\begin{equation}\label{eq:wrong}
 \sup_l\|e^{+B\widehat{\cL}(l)^a}\widehat v(l)\|_{\HS}<\infty.
\end{equation}
The positive sign in \eqref{eq:wrong} is the issue.  The source PDF included
with this packet contains the definitions, the inclusion of ordinary
distributions, and the conjecture on pages 147--149.

\section{Proof intuition}

A dual of a Gevrey test-function space must allow Fourier growth; a positive
exponential instead demands Fourier decay.  The simplest diagnostic is the
Dirac mass.  It is unquestionably in the target dual, while every one of its
Fourier blocks is the identity.  The least sub-Laplacian eigenvalue in the
spin-\(l\) block already tends to infinity, so even the two extremal diagonal
entries make the proposed norm explode.

For the repair, weighted \(L^2\) domains of \(e^{D\cL^a}\) form a Hilbert
scale.  Taking the dual reverses the weight, and exchanging Roumieu ``some
\(D\)'' with Beurling ``every \(D\)'' reverses the union/intersection
quantifier.  A small extra negative exponent converts uniform block control
to Peter--Weyl square summability.

\section{The counterexample}

\begin{theorem}\label{thm:counterexample}
For every \(1\le s<\infty\), the Dirac mass \(\delta_e\) at the identity of
\(\SU\) belongs to \((\gamma^{(s)}_{\mathbf X,L^2}(\SU))'\), but it does not
satisfy \eqref{eq:wrong} for any \(B>0\).  Consequently the printed second
clause of Conjecture 4.5.5 is false.
\end{theorem}

\begin{proof}
Source Definition 4.5.4 explicitly records
\[
 \mathcal D'(M)\subset
 (\gamma^s_{\mathbf X,L^2})'(M)\subset
 (\gamma^{(s)}_{\mathbf X,L^2})'(M).
\]
Thus \(\delta_e\) belongs to the asserted Beurling dual.  This also follows
directly from Sobolev point evaluation on the compact three-manifold \(\SU\).

For every irreducible unitary representation \(\pi_l\),
\[
 \widehat{\delta_e}(l)=\pi_l(e)^*=I_{d_l}.
\]
The smallest eigenvalue of \(\widehat{\cL}(l)\) is attained at
\(m=\pm l\) and equals \(l\).  Hence, for \(l>0\),
\begin{align*}
 \|e^{B\widehat{\cL}(l)^a}\widehat{\delta_e}(l)\|_{\HS}^2
 &=\sum_{m=-l}^{l}
       e^{2B(l(l+1)-m^2)^a}\\
 &\ge 2e^{2B l^a}.
\end{align*}
For every \(B>0\) this tends to infinity with \(l\), contradicting
\eqref{eq:wrong}.
\end{proof}

The failure is not marginal: every finite linear combination of derivatives
of \(\delta_e\) has at most polynomially growing Fourier blocks and is still
excluded by the positive exponential condition unless it vanishes.

\section{The corrected theorem}

We now make the natural topology explicit.  For \(D>0\), let
\[
 H_D=\operatorname{Dom}(e^{D\cL^a}),\qquad
 \|f\|_D=\|e^{D\cL^a}f\|_{L^2}.
\]
The source's \(\SU\) characterization identifies the Roumieu space with the
inductive union \(\bigcup_{D>0}H_D\).  Its stated Beurling analogue is the
projective intersection \(\bigcap_{D>0}H_D\).  We equip these with their
standard inductive- and projective-limit locally convex topologies.

\begin{theorem}[Exact sign correction]\label{thm:corrected}
Under these natural topologies:
\begin{enumerate}
 \item \(u\in(\gamma^s_{\mathbf X,L^2}(\SU))'\) if and only if
 \eqref{eq:roumieu-source} holds for every \(B>0\).
 \item \(v\in(\gamma^{(s)}_{\mathbf X,L^2}(\SU))'\) if and only if there
 exist \(B,K>0\) such that
 \begin{equation}\label{eq:correct}
  \|e^{-B\widehat{\cL}(l)^a}\widehat v(l)\|_{\HS}\le K
  \quad\text{for every }l\in\tfrac12\N_0.
 \end{equation}
\end{enumerate}
Thus the first printed clause is correct, and the second becomes correct
after changing \(+B\) to \(-B\).
\end{theorem}

\begin{proof}
Duality of a directed Hilbert scale gives
\begin{equation}\label{eq:dual-scales}
 \left(\bigcup_{D>0}H_D\right)'=\bigcap_{D>0}H_{-D},
 \qquad
 \left(\bigcap_{D>0}H_D\right)'=\bigcup_{D>0}H_{-D}.
\end{equation}
For completeness, the first identity follows because a functional on the
inductive limit is continuous on every \(H_D\).  For the second, continuity
is bounded by finitely many defining seminorms, hence by one of the directed
norms \(\|\cdot\|_D\).  Riesz representation on \(H_D\) then identifies the
corresponding dual with \(H_{-D}\).

Peter--Weyl and the spectral calculus give
\begin{equation}\label{eq:weighted-pw}
 \|w\|_{-D}^2
 =\sum_{l\in\frac12\N_0}d_l
   \|e^{-D\widehat{\cL}(l)^a}\widehat w(l)\|_{\HS}^2.
\end{equation}
Thus \(w\in H_{-D}\) immediately implies a uniform block bound with
exponent \(D\).

Conversely, suppose the block bound holds with exponent \(B\).  If \(D>B\),
then \eqref{eq:symbol} and \(\min_m(l(l+1)-m^2)=l\) for \(l>0\) yield
\begin{align*}
 \|w\|_{-D}^2
 &\le C_0+K^2\sum_{l>0}(2l+1)e^{-2(D-B)l^a}<\infty.
\end{align*}
The last series converges because \(a>0\); stretched-exponential decay
dominates the polynomial factor.  Therefore a uniform bound for every
\(B>0\) is equivalent to membership in every \(H_{-D}\), while a bound for
some \(B>0\) is equivalent to membership in some \(H_{-D}\).  Combine this
with \eqref{eq:dual-scales}.
\end{proof}

\section{Verification, novelty, and scope}

An exact-arithmetic checker verifies the block spectrum and extremal
eigenvalues for \(50{,}000\) half-integer representations, then stress-tests
the divergent lower bound and the corrected-weight tail over a grid of
\(s,B,D\).  These checks guard against indexing and sign mistakes; the proof
above establishes the infinite statements.

Four run indexes and bounded exact-title/core-keyword searches found no later
resolution through 11 August 2026.  The argument fully refutes the printed
Beurling clause and repairs it.  It does not address the distinct Heisenberg
conjecture, whose criterion is an \(L^1\)-integrability statement rather than
a uniform compact-group block bound.  Human review should focus on the
locally convex topology intended by the source for the strengthened corrected
theorem; Theorem~\ref{thm:counterexample} itself is topology-robust because
the source explicitly includes \(\mathcal D'(\SU)\).

\begin{thebibliography}{9}
\raggedright
\bibitem{source}
C. A. Taranto,
\emph{Wave equations on graded groups and hypoelliptic Gevrey spaces},
PhD thesis, Imperial College London; arXiv:1804.03544 (2018),
Corollary 4.3.6 and Section 4.5.1.
\end{thebibliography}

\end{document}
